\chapter{Microarchitecture du processeur}\label{ch:ll:cpu-microarchitecture}

En juin 2012, un programmeur demande pourquoi douze lignes de C++, additionnant les éléments d'un tableau supérieurs ou égaux à 128,
s'exécutent en 11.5~s sur des données aléatoires et en 1.9~s une fois ces mêmes données triées. Rien dans le source n'a
changé~; seul l'ordre des valeurs a changé, et avec lui la fréquence à laquelle le processeur devine correctement le sens dans lequel
un \texttt{if} va se dérouler. La réponse est devenue l'une des plus lues de l'histoire du site. Le chemin critique (hot path)
d'un système de trading ne compte que quelques milliers d'instructions, et le processeur qui les exécute n'est pas la machine séquentielle
que décrit le code source~: il recouvre des dizaines d'instructions, les exécute dans le désordre, devine les branchements et modifie
lui-même sa vitesse d'horloge. Ce chapitre décrit cette machine du point de vue du programme, mesure chaque effet sur l'ordinateur portable
de ce livre, et construit le banc d'essai que chaque chapitre ultérieur utilisera pour se chronométrer.

\section{Pipelines et exécution superscalaire}

\begin{definition}[Pipeline d'instructions, instructions par cycle]\label{def:ll:cpu-microarchitecture:pipeline}
Un \emph{pipeline d'instructions}\index{pipeline d'instructions} découpe l'exécution d'une instruction en étages (chargement,
décodage, renommage, ordonnancement, exécution, validation), chacun assuré par un matériel dédié, de sorte que de nombreuses instructions
sont en cours de traitement à la fois, chacune à un étage différent. Un pipeline superscalaire traite plusieurs instructions par étage et par cycle.
Les \emph{instructions par cycle}\index{instructions par cycle} (IPC) désignent le nombre d'instructions qu'un cœur achève en un cycle
de son horloge, moyenné sur un morceau de code.
\end{definition}

Un cœur de performance actuel décode et renomme six \omterm{def:ll:cpu-microarchitecture:pipeline}{instructions par cycle} et peut en valider six par cycle en régime soutenu,
lorsque les instructions sont indépendantes. Le code atteint rarement ce niveau~: un IPC compris entre un et trois est courant.
L'écart entre la largeur de la machine et l'IPC effectivement obtenu est le terrain de la conception orientée latence~: chaque
cycle au cours duquel moins d'instructions sont validées qu'il n'aurait été possible est un cycle perdu à attendre un résultat, la mémoire,
ou la reprise du pipeline après une mauvaise supposition.

\begin{omfigure}
\begin{tikzpicture}[font=\footnotesize,
  st/.style={draw=omDef,fill=omDef!8,rounded corners=2pt,minimum height=0.8cm,minimum width=1.55cm,align=center}]
\node[st] (f) at (0,0) {chargement\\ (prédiction)};
\node[st] (d) at (2.1,0) {décodage};
\node[st] (r) at (4.2,0) {renommage};
\node[st] (s) at (6.75,0) {ordonnancement};
\node[st,minimum height=2.0cm] (x) at (9.05,0) {exécution\\ (ALU,\\ lectures,\\ écritures)};
\node[st] (t) at (11.25,0) {validation};
\foreach \a/\b in {f/d,d/r,r/s,s/x,x/t} \draw[-Stealth,thick] (\a) -- (\b);
\draw[draw=omThm,fill=omThm!8,rounded corners=2pt] (3.4,-2.2) rectangle (12.05,-1.3);
\node[font=\footnotesize,align=center] at (7.725,-1.75) {tampon de réordonnancement~:\\ ordre du programme, achèvement dans le désordre};
\draw[-Stealth,omThm] (r.south) -- (4.2,-1.3);
\draw[-Stealth,omThm] (11.25,-1.3) -- (t.south);
\draw[-Stealth,thick,omProp] (x.north) to[out=120,in=60,looseness=0.6] node[pos=0.5,above,font=\footnotesize]{branchement mal prédit~: vidage et rechargement} (f.north);
\end{tikzpicture}

\omcaption{Pipeline d'un cœur à \omterm{def:ll:cpu-microarchitecture:ooo}{exécution dans le désordre}, simplifié. Les instructions entrent et sortent dans l'ordre du programme~; entre les deux,
l'ordonnanceur lance chacune dès que ses entrées sont prêtes, et le \omterm{def:ll:cpu-microarchitecture:ooo}{tampon de réordonnancement} conserve l'ordre du programme
pour la validation. Un branchement mal prédit écarte toutes les instructions chargées après lui.}\label{fig:ll:cpu-microarchitecture:pipeline}
\end{omfigure}

\section{Exécution dans le désordre et tampon de réordonnancement}

\begin{definition}[Exécution dans le désordre, tampon de réordonnancement]\label{def:ll:cpu-microarchitecture:ooo}
Dans l'\emph{exécution dans le désordre}\index{exécution dans le désordre} (out-of-order), le cœur lance chaque instruction dès que ses entrées
sont disponibles, et non dans l'ordre du programme, mais rend les résultats visibles dans l'ordre du programme. Le
\emph{tampon de réordonnancement}\index{tampon de réordonnancement} (reorder buffer) contient les instructions comprises entre le renommage et la validation~; sa taille borne
la distance à laquelle le cœur peut chercher du travail indépendant en avant d'une instruction bloquée.
\end{definition}

Le \omterm{def:ll:cpu-microarchitecture:ooo}{tampon de réordonnancement} de la famille de cœurs de performance de cet ordinateur portable contient 512 instructions. Lorsqu'un chargement échoue
dans tous les caches et attend cent nanosecondes la mémoire (chapitre~3), le cœur continue d'exécuter les instructions indépendantes qui
le suivent jusqu'à ce que le tampon soit plein, puis il se bloque. Ce qu'il ne peut pas faire, c'est raccourcir une chaîne de dépendances~: chaque
instruction qui a besoin du résultat de la précédente attend ce résultat, quelle que soit la largeur de la machine.

Pour observer ces effets, le livre chronomètre de petits morceaux de code isolés.

\begin{definition}[Microbenchmark, barrière de compilation]\label{def:ll:cpu-microarchitecture:ubench}
Un \emph{microbenchmark}\index{microbenchmark} chronomètre un petit morceau de code isolé, de nombreuses fois, afin de mesurer une
propriété de la machine ou du code. Une \emph{barrière de compilation}\index{barrière de compilation} est une instruction que le compilateur
doit traiter comme une lecture ou une écriture en mémoire (ou sur une valeur donnée), bien qu'elle n'émette aucune instruction~; elle empêche
l'optimiseur de supprimer ou de déplacer le code chronométré.
\end{definition}

Un \omterm{def:ll:cpu-microarchitecture:ubench}{microbenchmark} ne mesure que ce qu'il mesure~: une boucle sur des données qui restent dans le cache de premier niveau, sur un cœur
déjà chauffé, une fois ses branchements appris. C'est le meilleur cas du chemin critique, et non son comportement après une période calme
(caches froids, cœur dans un état de faible consommation), qui est souvent ce qui compte en trading (chapitre~5).

\begin{proposition}[Chaînes nécessaires pour saturer les unités]\label{prop:ll:cpu-microarchitecture:little}
Si une opération a une latence de $\ell$ cycles et que le cœur peut en lancer $w$ par cycle, une boucle a besoin d'au moins $\ell w$
chaînes de dépendances indépendantes pour en lancer $w$ par cycle~; avec $k < \ell w$ chaînes, elle achève au plus $k/\ell$ opérations par cycle.
\end{proposition}

\begin{proof}
Chaque chaîne a au plus une opération en cours, qui occupe $\ell$ cycles~; $k$ chaînes ont donc au plus $k$ opérations en cours. Par
la loi de Little, le débit d'achèvement est au plus $k/\ell$, qui n'atteint $w$ que lorsque $k \ge \ell w$.
\end{proof}

La \Cref{fig:ll:cpu-microarchitecture:chains} somme 4\,096 doubles, conservés dans le cache, avec de un à seize accumulateurs. Une
addition de flottants sur ce cœur a une latence de deux cycles lorsqu'elle est suivie d'une autre addition, et deux peuvent être lancées
par cycle~: la \cref{prop:ll:cpu-microarchitecture:little} prédit quatre chaînes, et le coût mesuré par addition est divisé par deux
de un à deux accumulateurs, de nouveau divisé par deux à quatre, puis ne s'améliore plus au-delà. Le compilateur ne le fait pas de lui-même~:
sans autorisation de réassocier les additions de flottants (\texttt{-ffast-math}, chapitre~8), il doit conserver la seule
chaîne écrite dans le source.

\begin{omfigure}
\begin{tikzpicture}
\begin{axis}[width=9cm,height=5.2cm,ybar,bar width=14pt,ymin=0,enlarge y limits={upper,value=0.2},
  symbolic x coords={1,2,4,8,16},xtick=data,enlarge x limits=0.15,
  xlabel={accumulateurs indépendants},ylabel={ns par addition},tick label style={font=\small},label style={font=\small},
  grid=major,grid style={gray!25},nodes near coords,nodes near coords style={font=\footnotesize},
  point meta=y,every node near coord/.append style={/pgf/number format/fixed,/pgf/number format/precision=2},
  table/col sep=comma]
\addplot[fill=omDef!55,draw=omDef] table[x=accumulators,y=ns_per_add]{figdata/low-latency/02-cpu-microarchitecture/measured_chains.csv};
\end{axis}
\end{tikzpicture}

\omcaption{Coût par addition de flottants lors de la somme de 4\,096 doubles avec 1 à 16 accumulateurs indépendants. Deux
cycles de latence et deux additionneurs suffisent à atteindre quatre chaînes (\cref{prop:ll:cpu-microarchitecture:little}). Mesuré sur un
ordinateur portable (Intel Core Ultra 7 155H) sous WSL2, sans cœurs isolés, épinglé sur un seul CPU, meilleur résultat sur 30 exécutions. Données~:
\texttt{bench\_cpu.py}.}\label{fig:ll:cpu-microarchitecture:chains}
\end{omfigure}

Le cœur peut aussi supprimer les dépendances que contient le programme. Une chaîne de dix additions dépendantes d'un registre s'exécute à
une addition par cycle, comme il se doit~; la même chaîne ajoutant la constante 1 s'exécutait à environ cinq additions par
cycle sur cet ordinateur portable, parce que le renommeur de cette famille de cœurs résout lui-même les additions à petite immédiate, jusqu'à six par cycle, sans
les envoyer à une unité d'exécution. Un \omterm{def:ll:cpu-microarchitecture:ubench}{microbenchmark} qui utilise \texttt{add \$1} pour compter des cycles mesure le renommeur~;
le harnais ci-dessous mesure l'horloge avec la forme à registre.

\section{Prédiction de branchement}

\begin{definition}[Prédiction de branchement, erreur de prédiction de branchement]\label{def:ll:cpu-microarchitecture:branch}
La \emph{prédiction de branchement}\index{prédiction de branchement} (branch prediction) est la supposition du cœur, formulée à la lecture, sur la direction et la cible de
chaque branchement, afin que le pipeline puisse continuer à lire les instructions avant que le branchement soit exécuté. Une \emph{erreur de
prédiction de branchement}\index{erreur de prédiction de branchement} (branch misprediction) est une supposition fausse~: le cœur jette chaque instruction lue après le
branchement et relit depuis la bonne cible, ce qui fait perdre autant de cycles que la profondeur de l'avant du pipeline.
\end{definition}

\begin{proposition}[Quand le code sans branchement gagne]\label{prop:ll:cpu-microarchitecture:breakeven}
Un corps de boucle comportant un branchement dépendant des données coûte en moyenne $c_0 + mP$ cycles par itération, où $m$ est le
taux d'erreurs de prédiction et $P$ la pénalité d'erreur de prédiction~; une réécriture sans branchement qui fait toujours le travail des deux côtés,
coûte $c_0 + c$. La réécriture est moins chère exactement lorsque $m > c/P$.
\end{proposition}

\begin{proof}
Chaque itération subit une erreur de prédiction avec probabilité $m$ et paie alors $P$~; la linéarité de l'espérance donne $c_0 + mP$. Comparer
à $c_0 + c$.
\end{proof}

La \Cref{fig:ll:cpu-microarchitecture:branch} reprend la boucle de 2012. Sur des données triées, le prédicteur a raison à chaque itération sauf
une sur 32\,768 itérations, et la boucle avec branchements est la plus rapide~; sur des données mélangées, il a raison une fois sur deux, et la boucle
est environ treize fois plus lente. La boucle sans branchement, qui masque chaque élément par des uns ou des zéros, coûte autant
dans les deux cas. La moitié des itérations subissent une erreur de prédiction sur des données mélangées, donc l'écart entre les deux barres avec branchements est la
moitié de la pénalité par élément~: environ \qty{5.2}{\nano\second}, soit quelque 24 cycles à l'horloge mesurée de ce cœur, un peu au-dessus de la
valeur «~parfois plus de 20 cycles d'horloge~» qu'une mesure publiée donne pour cette famille de cœurs. La boucle sans branchement
coûte environ \qty{0.11}{\nano\second} de plus par élément sur des données triées, donc, selon la
\cref{prop:ll:cpu-microarchitecture:breakeven}, elle l'emporte dès que plus d'environ deux itérations sur cent
subissent une erreur de prédiction.

\begin{omfigure}
\begin{tikzpicture}
\begin{axis}[width=9cm,height=5.2cm,ybar,bar width=16pt,ymin=0,
  symbolic x coords={sorted,shuffled},xtick=data,xticklabels={données triées,données mélangées},enlarge x limits=0.5,
  ylabel={ns par élément},tick label style={font=\small},label style={font=\small},
  grid=major,grid style={gray!25},
  legend columns=2,legend style={font=\footnotesize,at={(0.5,-0.2)},anchor=north,draw=none,fill=none,
  /tikz/every even column/.append style={column sep=8pt}},table/col sep=comma]
\addplot[fill=omThm!45,draw=omThm] table[x=order,y=branchy]{figdata/low-latency/02-cpu-microarchitecture/measured_branch.csv};
\addplot[fill=omDef!45,draw=omDef] table[x=order,y=branchless]{figdata/low-latency/02-cpu-microarchitecture/measured_branch.csv};
\legend{avec branchement,sans branchement}
\end{axis}
\end{tikzpicture}

\omcaption{La boucle de la question de 2012 (sommer les éléments $\ge 128$ de 32\,768 octets aléatoires) avec un branchement et avec un
masque, sur des données triées et mélangées. Mesuré sur un ordinateur portable (Intel Core Ultra 7 155H) sous WSL2, sans cœurs isolés,
épinglé sur un seul CPU, meilleur résultat sur 30 exécutions. Données~: \texttt{bench\_cpu.py}.}\label{fig:ll:cpu-microarchitecture:branch}
\end{omfigure}

La leçon pour le trading n'est pas «~supprimer les branchements~». La plupart des branchements sur le chemin critique (hot path) sont prévisibles (le type de message est
généralement le même, le contrôle de risque passe le plus souvent) et ne coûtent presque rien. Les plus coûteux dépendent du marché~:
de quel côté se situe un message, si un prix améliore le carnet d'ordres. Ils subissent des erreurs de prédiction par rafales, précisément quand le marché est
animé, et il vaut la peine de les mesurer un à un.

\section{Variation dynamique de fréquence, turbo et états de veille}

\begin{definition}[Variation dynamique de fréquence, état de veille]\label{def:ll:cpu-microarchitecture:dvfs}
La \emph{variation dynamique de fréquence}\index{variation dynamique de fréquence} (dynamic frequency scaling) modifie l'horloge (et la tension) d'un cœur pendant qu'il tourne,
selon la charge, la température, le budget de puissance et le nombre de cœurs actifs~; les fréquences les plus élevées, appelées turbo,
ne sont disponibles que tant que le boîtier reste dans ses limites. Un \emph{état de veille}\index{etat@état de veille} (C-state) est un
état à faible consommation qu'un cœur adopte lorsqu'il n'a rien à exécuter, et dont le réveil prend quelques microsecondes, les états
plus profonds demandant davantage de temps.
\end{definition}

\begin{definition}[Compteur d'horodatage]\label{def:ll:cpu-microarchitecture:tsc}
Le \emph{compteur d'horodatage}\index{compteur d'horodatage} (TSC) est un registre de 64 bits de chaque cœur x86 qui compte à une
cadence constante sur les processeurs actuels, quelle que soit l'horloge du cœur, lu par les instructions \texttt{rdtsc} et
\texttt{rdtscp} en quelques nanosecondes.
\end{definition}

Les deux définitions interagissent. Le TSC compte du temps, non des cycles~: sur cet ordinateur portable, il cadence à environ 3.0~GHz pendant que le cœur
fonctionne à une fréquence qui va de son horloge de veille à plusieurs gigahertz. Un \omterm{def:ll:cpu-microarchitecture:ubench}{microbenchmark} qui affiche les battements du TSC comme des «~cycles~» se trompe d'un
facteur égal à ce rapport. Le harnais affiche des nanosecondes et, lorsque les cycles comptent, mesure l'horloge du cœur avec une chaîne d'additions dépendantes
sur des registres, un cycle chacune. Sur un serveur, la leçon est opérationnelle~: un cœur inactif depuis un certain temps
répond à son premier message depuis un \omterm{def:ll:cpu-microarchitecture:dvfs}{état de veille} profond et à une horloge basse, si bien qu'un hôte de trading limite les \omterm{def:ll:cpu-microarchitecture:dvfs}{états de veille} et fixe
sa fréquence (chapitre~13).

\begin{dated}{2026-09}{L'ordinateur portable utilisé pour les mesures de ce livre}\label{dat:ll:cpu-microarchitecture:laptop}
L'Intel Core Ultra 7 155H, lancé au T4 2023, compte 16 cœurs~: 6 cœurs de performance dotés chacun de deux
threads matériels, 8 cœurs efficaces et 2 cœurs efficaces basse consommation, 22 threads au total, avec des fréquences turbo maximales de
4.8~GHz sur les cœurs de performance et 3.8~GHz sur les cœurs efficaces, et 24~Mo de cache partagé (page
de spécification d'Intel).
\end{dated}

\section{Multithreading simultané et cœurs hétérogènes}

\begin{definition}[Multithreading simultané, cœurs hétérogènes]\label{def:ll:cpu-microarchitecture:smt}
Dans le \emph{multithreading simultané}\index{multithreading simultané} (SMT, hyper-threading), un cœur exécute deux
threads matériels, qui partagent ses unités d'exécution, ses caches et ses prédicteurs~; chacun apparaît au système d'exploitation comme un CPU.
Un processeur possède des \emph{cœurs hétérogènes}\index{cœurs hétérogènes} (heterogeneous cores) lorsque ses cœurs diffèrent par leur conception, typiquement de grands
cœurs de performance et de petits cœurs efficaces, aux vitesses différentes pour un même code.
\end{definition}

Les deux comptent pour un thread sensible à la latence. Son thread matériel jumeau se dispute les mêmes unités d'exécution et
les caches de premier niveau, si bien que les hôtes de trading désactivent soit le SMT, soit laissent inactif le jumeau de chaque cœur critique. Sur un
processeur hétérogène, le même code s'exécute à des vitesses différentes selon l'endroit où le système d'exploitation le place~; un
thread critique doit être épinglé sur un cœur de performance.

Cet ordinateur portable ajoute une couche que n'ont généralement pas les serveurs installés dans un site de colocation~: il fait tourner Linux dans une
machine virtuelle. Les 22 CPU que voit Linux sont des CPU virtuels, que l'hyperviseur planifie sur le cœur physique de son choix,
et le compte rendu du processeur sur les types de ses cœurs est masqué à l'invité. La \Cref{fig:ll:cpu-microarchitecture:freq}
mesure l'horloge perçue par chaque CPU virtuel~: les cœurs de performance et les cœurs efficaces n'apparaissent pas en deux groupes, et
épingler un thread sur un CPU virtuel n'équivaut pas à l'épingler sur un seul cœur physique. Les mesures de ce livre sont donc les
valeurs au mieux de l'ordinateur portable, et chaque légende le signale.

\begin{omfigure}
\begin{tikzpicture}
\begin{axis}[width=12cm,height=5cm,ybar,bar width=6pt,ymin=0,ymax=5.5,xmin=-0.8,xmax=21.8,
  xtick={0,2,4,6,8,10,12,14,16,18,20},xlabel={CPU virtuel},ylabel={horloge mesurée (GHz)},
  tick label style={font=\small},label style={font=\small},grid=major,grid style={gray!25},table/col sep=comma]
\addplot[fill=omDef!55,draw=omDef] table[x=cpu,y=ghz_reg]{figdata/low-latency/02-cpu-microarchitecture/measured_freq.csv};
\end{axis}
\end{tikzpicture}

\omcaption{L'horloge à laquelle chacun des 22 CPU virtuels de l'ordinateur portable exécute une chaîne d'additions dépendantes dans des registres (une
addition par cycle), meilleure valeur sur 15 exécutions. L'hyperviseur masque quel cœur physique exécute quel CPU virtuel~; les six
cœurs de performance et les dix cœurs efficaces n'apparaissent pas en groupes. Mesuré sur un ordinateur portable (Intel Core Ultra 7 155H) sous WSL2,
sans cœurs isolés, la machine étant par ailleurs au repos. Données~: \texttt{bench\_cpu.py}.}\label{fig:ll:cpu-microarchitecture:freq}
\end{omfigure}

\section{Tutoriel~: chronométrer honnêtement un noyau}\label{tut:ll:cpu-microarchitecture:tut}

\begin{tutorial}
\textbf{Objectif.} Construire le harnais de \omterm{def:ll:cpu-microarchitecture:ubench}{microbenchmark}, puis l'utiliser pour mesurer une \omterm{def:ll:cpu-microarchitecture:branch}{erreur de prédiction de branchement}, une chaîne de dépendances et
l'horloge. \textbf{État final~:}
\cref{fig:ll:cpu-microarchitecture:branch,fig:ll:cpu-microarchitecture:chains,fig:ll:cpu-microarchitecture:freq} et une
estimation de la pénalité de mauvaise prédiction en cycles.
\begin{tutsteps}
\item \textbf{Garder le code vivant et encadrer le compteur.} Une instruction \texttt{asm} vide qui prend la valeur en entrée est une
\omterm{def:ll:cpu-microarchitecture:ubench}{barrière de compilation}~: l'optimiseur doit calculer la valeur et ne peut pas déplacer de travail au-delà. La région chronométrée
commence par \texttt{lfence; rdtsc} et se termine par \texttt{rdtscp; lfence}, afin que le travail antérieur soit achevé avant la
première lecture et que la seconde lecture attende la fin de la région.
\omcode{code/firm/ubench/cpp/firm_ubench.hpp}{22}{41}{Le harnais C++20~: barrières et lectures encadrées du compteur d'horodatage.}\label{lst:ll:cpu-microarchitecture:cpp}
\item \textbf{Le même en Rust.} Rust dispose de \texttt{std::hint::black\_box} pour la barrière~; le compteur demande les
intrinsèques de \texttt{std::arch} et un bloc \texttt{unsafe}, dont l'argument de sûreté figure dans un commentaire.
\omcode{code/firm/ubench/rust/src/lib.rs}{7}{27}{Les lectures encadrées de la réplique Rust.}\label{lst:ll:cpu-microarchitecture:rust}
\item \textbf{Les noyaux.} La boucle à branchement conserve son branchement parce qu'une instruction \texttt{asm} opaque s'y trouve~;
sans elle, le compilateur peut transformer le \texttt{if} en mouvement conditionnel et l'expérience disparaît.
\omcode{code/low-latency/02-cpu-microarchitecture/cpp/ll_cpu.hpp}{9}{31}{La boucle de 2012 avec un branchement et avec un masque.}\label{lst:ll:cpu-microarchitecture:kernels}
\item \textbf{Lancer} \texttt{python bench\_cpu.py}~: il compile \texttt{ll\_cpu\_bench.cpp}, exécute les trois modes épinglés
avec \texttt{taskset}, et écrit les CSV mesurés ainsi que leurs fichiers \texttt{.meta}~; \texttt{ll\_cpu.penalty()} transforme
les valeurs de branchement en une pénalité en cycles.
\end{tutsteps}
\textbf{Pour aller plus loin.} Retirer l'instruction \texttt{asm} de \texttt{sum\_branchy}, examiner le désassemblage
(\texttt{objdump -d}) et mesurer de nouveau~; remplacer les octets aléatoires par un motif de période 8 et observer le prédicteur
l'apprendre.
\end{tutorial}

\section{Construction~: le harnais de microbenchmark}\label{bld:ll:cpu-microarchitecture:ubench}

\begin{build}
\bfield{Rôle} Une façon commune de mesurer le temps de code dans tout le livre, en C++20, en Rust et dans les pilotes Python qui produisent chaque
figure mesurée, afin que les chiffres de chapitres différents soient comparables et indiquent leur provenance.
\bfield{Interface} C++20 \texttt{firm::ubench}~: \texttt{do\_not\_optimize(x)}, \texttt{clobber()},
\texttt{tsc\_start()}, \texttt{rdtscp()}, \texttt{Clock::calibrate(ms)} et \texttt{Clock::ns(ticks)},
\texttt{overhead()}, \texttt{run(fn, reps, warm)}, \texttt{quantile(v, p)}. \texttt{firm\_ubench} en Rust~: les mêmes
fonctions. \texttt{firm\_ubench} en Python~: \texttt{compile\_cpp}, \texttt{run(exe, *args, cpus)},
\texttt{machine()}, \texttt{write\_measured(csv, header, rows, **meta)}, \texttt{quantiles}.
\bfield{Règles} Chaque région chronométrée est encadrée~; les résultats sont en nanosecondes (les ticks n'apparaissent que dans le harnais)~; chaque CSV mesuré
est accompagné d'un fichier \texttt{.meta} qui nomme la machine, le compilateur, les options, la liste des CPU et la charge~; les pilotes s'appellent
\texttt{bench\_*.py} et ne sont jamais lancés par \texttt{make figdata}.
\bfield{Tests d'acceptation} \texttt{code/firm/ubench/}~: calibration entre 0.5 et 6 ticks par nanoseconde, un compteur monotone, un travail
plus coûteux qu'une région vide (C++, Rust), et l'aller-retour du CSV et des métadonnées du pilote (Python).
\bfield{Pour aller plus loin} Lire le compteur de cycles du cœur par \texttt{perf\_event\_open} lorsque cela est permis~; indiquer l'étendue des
exécutions répétées avec chaque chiffre.
\end{build}

\begin{omsources}
\item Stack Overflow, ``Why is processing a sorted array faster than processing an unsorted array?'', question
11227809, 2012.
\item A.~Fog, \emph{The microarchitecture of Intel, AMD and VIA CPUs}, mis à jour en 2026.
\item Chips and Cheese, ``Popping the hood on Golden Cove'' (2021) et ``Lion Cove: Intel's P-core roars'' (2024).
\item Intel, \emph{Intel 64 and IA-32 Architectures Optimization Reference Manual}.
\end{omsources}

\section{Exercices}

\begin{exercise}[$\star$]\label{exo:ll:cpu-microarchitecture:1}
Un cœur retire 2.4 \omterm{def:ll:cpu-microarchitecture:pipeline}{instructions par cycle} à 4.0~GHz. Combien d'instructions un chemin critique de \qty{600}{\nano\second}
exécute-t-il~?
\end{exercise}

\begin{exercise}[$\star$]\label{exo:ll:cpu-microarchitecture:2}
Un \omterm{def:ll:cpu-microarchitecture:ubench}{microbenchmark} indique 3\,000 ticks TSC pour une fonction sur un ordinateur portable dont le TSC cadence à 3.0~GHz tandis que le cœur tourne à
4.5~GHz. Combien de temps la fonction a-t-elle pris, et combien de cycles de cœur~?
\end{exercise}

\begin{exercise}[$\star$]\label{exo:ll:cpu-microarchitecture:3}
Une opération a une latence de 4 cycles et le cœur peut en lancer deux par cycle. Combien de chaînes indépendantes une boucle
doit-elle comporter pour atteindre le débit du cœur, et quelle fraction de ce débit atteint une boucle à trois chaînes~?
\end{exercise}

\begin{exercise}[$\star\star$]\label{exo:ll:cpu-microarchitecture:4}
Un gestionnaire de messages comporte un branchement, selon le côté du message, qui se trompe dans 30~\% des cas, avec une pénalité de 20 cycles. Une
version sans branchement coûte 4 cycles supplémentaires à chaque message. Laquelle est la plus rapide, et de combien par message~?
\end{exercise}

\begin{exercise}[$\star\star$]\label{exo:ll:cpu-microarchitecture:5}
À partir de la \cref{fig:ll:cpu-microarchitecture:branch}, estimer la pénalité de mauvaise prédiction en nanosecondes, puis en cycles à
l'horloge de la \cref{fig:ll:cpu-microarchitecture:freq} pour le CPU~2. Sur quelle hypothèse concernant les données mélangées repose
l'estimation~?
\end{exercise}

\begin{exercise}[$\star\star$]\label{exo:ll:cpu-microarchitecture:6}
Pourquoi le \omterm{def:ll:cpu-microarchitecture:ooo}{tampon de réordonnancement} ne peut-il pas masquer une chaîne de dépendances, quelle que soit sa taille~? Que masque-t-il à la place~?
\end{exercise}

\begin{exercise}[$\star\star\star$]\label{exo:ll:cpu-microarchitecture:7}
\emph{Programmation.} Ajouter un mode à \texttt{ll\_cpu\_bench.cpp} qui exécute la boucle à branchement sur des données dont les valeurs alternent par
blocs de $b$ éléments au-dessus et au-dessous de 128, pour $b = 1, 2, 4, 8, 64$. Qu'apprend le prédicteur, et pour quel $b$ la boucle atteint-elle la
vitesse du cas trié~?
\end{exercise}

\begin{exercise}[$\star\star\star$]\label{exo:ll:cpu-microarchitecture:8}
\emph{Trouver la faille.} «~Nous avons mesuré notre encodeur d'ordres à 40 cycles avec une boucle d'instructions \texttt{add \$1} comme
référence de cycle, sur un cœur déjà chaud, en l'appelant dix millions de fois de suite sur le même ordre.~»
\end{exercise}

\section{Problème~: Le branchement qui a coûté une cotation}

\begin{problem}[{Problème du week-end --- un \texttt{if} dans le chemin de cotation}]\label{pb:ll:cpu-microarchitecture:1}
Le chemin critique du moteur de cotation comporte un seul branchement dépendant des données~: savoir si une transaction entrante s'est faite au cours acheteur ou au cours vendeur.
Un jour calme, les deux côtés alternent par longues séries~; un jour chargé, ils sont presque aléatoires. L'équipe mesure la boucle de 2012 comme
substitut (la \cref{fig:ll:cpu-microarchitecture:branch}) et doit décider s'il faut réécrire ce branchement.

\medskip\noindent\textbf{Partie I --- La mesure.}
\begin{enumerate}
\item Quels sont les quatre coûts mesurés par élément~?
\item Pourquoi la boucle à branchements sur les données triées est-elle plus rapide que la boucle sans branchement~?
\item Estimer la pénalité d'une \omterm{def:ll:cpu-microarchitecture:branch}{erreur de prédiction de branchement} en nanosecondes.
\item La convertir en cycles avec la fréquence d'horloge mesurée du CPU~2.
\end{enumerate}

\medskip\noindent\textbf{Partie II --- Le modèle.}
\begin{enumerate}[resume]
\item Écrire le coût attendu par itération des boucles à branchements et sans branchement en fonction du
taux d'\omterm{def:ll:cpu-microarchitecture:branch}{erreur de prédiction de branchement}.
\item À quel taux d'erreur de prédiction les deux versions coûtent-elles autant~?
\item Un jour calme, le branchement se trompe sur 1~\% des transactions. Quelle version l'emporte, et de combien par transaction~?
\item Un jour chargé, il se trompe sur 35~\%. Laquelle l'emporte, et de combien~?
\end{enumerate}

\medskip\noindent\textbf{Partie III --- Ce que cela vaut.}
\begin{enumerate}[resume]
\item Un jour chargé, les transactions arrivent par rafales, à 50\,000 par seconde. Combien de temps de traitement
par seconde la version à branchements gaspille-t-elle~?
\item Par rapport au budget de \qty{1.3}{\micro\second} du chemin critique pour les étapes logicielles (chapitre~1), la pénalité par
transaction mal prédite est-elle grande ou petite~?
\item Pourquoi les erreurs de prédiction se regroupent-elles les jours qui comptent~?
\item Que changerait un prédicteur qui aurait appris le motif du jour chargé~?
\end{enumerate}

\medskip\noindent\textbf{Partie IV --- La décision.}
\begin{enumerate}[resume]
\item Énoncer le \emph{résultat nommé}~: la pénalité d'erreur de prédiction mesurée sur cet ordinateur portable, en nanosecondes et en cycles, et le
taux d'erreur de prédiction au-delà duquel la réécriture sans branchement l'emporte.
\item Pourquoi la pénalité mesurée dépasse-t-elle la valeur publiée pour la famille de cœurs~?
\item Que mesure le TSC dans cette expérience, et pourquoi ne s'agit-il pas de cycles~?
\item Pourquoi épingler le benchmark sur un seul CPU, et que signifie l'épinglage sous un hyperviseur~?
\item Que ferait un compilateur de la boucle à branchements si l'instruction \texttt{asm} était supprimée~?
\item Citer un autre branchement du chemin critique de trading qui est prévisible, et un qui ne l'est pas.
\item Comment mesurer le taux d'erreur de prédiction sur l'hôte de production~?
\item En une phrase~: quand vaut-il la peine d'écrire du code sans branchement~?
\end{enumerate}
\end{problem}

\section{Questions d'entretien}

\begin{interviewq}[$\star$ \iqroles{développeur}]\label{iq:ll:cpu-microarchitecture:1}
Pourquoi trier un tableau peut-il rendre une boucle qui le parcourt plusieurs fois plus rapide, alors que le travail de la boucle ne dépend pas de l'ordre~?
\end{interviewq}

\begin{interviewq}[$\star\star$ \iqroles{développeur}]\label{iq:ll:cpu-microarchitecture:2}
Expliquez l'\omterm{def:ll:cpu-microarchitecture:ooo}{exécution dans le désordre} (out-of-order). Qu'est-ce qui limite la quantité de travail indépendant que le cœur peut trouver~?
\end{interviewq}

\begin{interviewq}[$\star\star$ \iqroles{développeur}]\label{iq:ll:cpu-microarchitecture:3}
Pourquoi la somme d'un tableau avec quatre accumulators s'exécute-t-elle plus vite qu'avec un seul~? Pourquoi le compilateur ne le fait-il pas pour vous~?
\end{interviewq}

\begin{interviewq}[$\star\star$ \iqroles{développeur}]\label{iq:ll:cpu-microarchitecture:4}
Vous voulez chronométrer une fonction qui prend 50 nanosecondes. Comment vous y prenez-vous, et que peut-il mal se passer~?
\end{interviewq}

\begin{interviewq}[$\star\star$ \iqroles{développeur}]\label{iq:ll:cpu-microarchitecture:5}
Un thread critique pour la latence doit-il partager un cœur avec son thread jumeau en hyper-threading~? Pourquoi~?
\end{interviewq}

\begin{interviewq}[$\star\star\star$ \iqroles{développeur}]\label{iq:ll:cpu-microarchitecture:6}
Le premier ordre après une minute calme est toujours plus lent que les suivants. Citez trois causes dans le processeur et
comment vous confirmeriez chacune.
\end{interviewq}
